\documentclass[11pt,a4paper]{article}
\usepackage[utf8]{inputenc}
\usepackage[T1]{fontenc}
\usepackage[margin=2.5cm]{geometry}
\usepackage{amsmath,amssymb,amsthm,booktabs,hyperref}
\newtheorem{theorem}{Theorem}[section]
\newtheorem{lemma}[theorem]{Lemma}
\newtheorem{proposition}[theorem]{Proposition}
\newcommand{\F}{\mathbb F_2}
\newcommand{\wt}{\operatorname{wt}}
\newcommand{\Fix}{\operatorname{Fix}}
\title{The Cubic Case of the Homogeneous Rotation-Symmetric Bent Function Conjecture}
\author{Author details to be completed}
\date{Revised working manuscript --- October 2026}
\begin{document}
\maketitle
\begin{abstract}
We prove that no homogeneous rotation-symmetric Boolean function of algebraic degree three is bent in any positive even dimension. Writing \(n=2^s m\) with \(m\) odd, an odd-order fixed-space argument reduces a hypothetical cubic homogeneous rotation-symmetric bent function to a rotation-symmetric bent function of degree at most three in \(2^s\) variables. Cubic orbit folding cannot produce the antipodal quadratic orbit, whereas an antipodal rigidity theorem proved here shows that every rotation-symmetric bent function of degree at most three in power-of-two dimension at least four must contain that orbit. The remaining two-variable case is immediate. Earlier exact computations in dimensions 16 and 32 are retained as independent validation, not as logical ingredients.
\end{abstract}
\noindent\textbf{Keywords:} Boolean functions; bent functions; rotation symmetry; homogeneous Boolean functions; derivatives; fixed spaces.

\section{Scope and relation to earlier work}
St\u{a}nic\u{a} and Maitra~\cite{stanica} conjectured that homogeneous rotation-symmetric bent functions of degree greater than two do not exist. Earlier partial results include those of Meng, Chen and Fu~\cite{meng}, Zhang and Gao~\cite{zhang}, and Cusick and Sanger~\cite{cusick}. Meng, Chen and Fu exclude several support-dependent classes; Cusick and Sanger obtain strong exclusions in dimensions \(n=2p\) with \(p\) an odd prime.

Sun, Shi, Liu and Fu~\cite{sun} give further nonexistence criteria, including results for cubic homogeneous rotation-symmetric functions based on balanced derivatives, SANF/orbit structure and arithmetic hypotheses on \(n\). Their work is particularly close in spirit to one step of the present argument: derivatives of cubic functions are quadratic. The proof here uses that observation together with an odd-order fixed-space reduction and a rigidity property in power-of-two dimension.

The main theorem below is uniform in the even dimension: no restriction on the odd part of \(n\) or on \(v_2(n)\) remains. The comparison with the cited papers is intended to delimit the mechanism and scope of the result; it is not an exhaustive priority claim. A separate bibliographic review is required before asserting that the theorem is the first result of its exact form.

The proof of the main theorem is algebraic. Exact computations previously obtained in dimensions \(16\) and \(32\), as well as independent folding and quadratic checks, are retained later as reproducibility controls and historical motivation, but they are not logical premises. Nonhomogeneity must not be conflated with homogeneity: rotation-symmetric cubic bent functions with lower-degree terms are known and are compatible with the antipodal obstruction proved below.

\section{Notation}
Let $V_n=\F^n$. The algebraic normal form (ANF) of a Boolean function is its unique multilinear polynomial over $\F$. Homogeneity of degree three means that every nonzero ANF term has degree three. The zero function may be included in spans without affecting any nonexistence assertion.
For $a\in V_n$, set $D_a f(x)=f(x+a)+f(x)$. Define
\[
 W_f(b)=\sum_{x\in V_n}(-1)^{f(x)+b\cdot x},\qquad
 W_f(0)=2^n-2\wt(f).
\]
A function is balanced if $W_f(0)=0$, and bent if $|W_f(b)|=2^{n/2}$ for every $b$. Equivalently, every nonzero directional derivative is balanced~\cite{rothaus}. Indeed, the Fourier transform of the autocorrelation $A_f(a)=\sum_x(-1)^{D_af(x)}$ is $W_f(b)^2$. Thus constant squared Walsh magnitude is equivalent to $A_f(0)=2^n$ and $A_f(a)=0$ for $a\ne0$.

Let $\sigma$ be cyclic rotation of the coordinates. Rotation symmetry means $f\circ\sigma=f$. Such functions are linear combinations of sums over distinct cyclic monomial orbits. All polynomial sums are in $\F$; exponential sums and weights are ordinary integers.

\section{Odd-order reduction}
\begin{lemma}\label{lem:reduction}
Let $V$ be a finite-dimensional vector space over $\F$, and let $T\in\operatorname{GL}(V)$ satisfy $T^m=I$ for an odd positive integer $m$. If $q:V\to\F$ has degree at most two and $q\circ T=q$, then $q$ is balanced on $V$ if and only if $q|_U$ is balanced on $U=\Fix(T)$.
\end{lemma}
\begin{proof}
Adding $q(0)$ preserves balance both on $V$ and on $U$, so assume $q(0)=0$. Its polar form
\[B(x,y)=q(x+y)+q(x)+q(y)\]
is alternating bilinear and $T$-invariant. Set $P=\sum_{j=0}^{m-1}T^j$. Then $TP=PT=P$, $Pu=u$ for $u\in U$, and $P^2=P$. Hence $V=U\oplus K$, where $K=\ker P$. For $u\in U$ and $k\in K$, invariance and oddness give
\[B(u,k)=\sum_{j=0}^{m-1}B(u,T^jk)=B(u,Pk)=0.\]
It follows that $q(u+k)=q(u)+q(k)$, and thus
\[S_V(q)=S_U(q|_U)S_K(q|_K),\qquad S_E(h)=\sum_{x\in E}(-1)^{h(x)}.\]
We show $S_K(q|_K)\ne0$. Put $R=\{z\in K:B(z,k)=0\text{ for all }k\in K\}$. Since $P$ commutes with $T$, $K$ is invariant under $T$ and $T^{-1}$. If $z\in R$ and $k\in K$, then $B(Tz,k)=B(z,T^{-1}k)=0$, so $R$ is $T$-invariant. On $R$, the function $q$ is linear. More explicitly, if $z\in R$, then all $T^jz$ lie in $R$, so every cross-polar term $B(T^iz,T^jz)$ vanishes. Hence
\[
q(Pz)=q\!\left(\sum_{j=0}^{m-1}T^jz\right)
     =\sum_{j=0}^{m-1}q(T^jz)
     =m q(z)=q(z).
\]
Since $z\in K$, $Pz=0$, and therefore $q(z)=q(0)=0$.
Character orthogonality now gives
\begin{align*}
 S_K(q|_K)^2
 &=\sum_{z\in K}(-1)^{q(z)}\sum_{x\in K}(-1)^{B(x,z)}\\
 &=|K|\sum_{z\in R}(-1)^{q(z)}=|K||R|>0.
\end{align*}
The factorization proves the equivalence of balance.
\end{proof}

\begin{proposition}\label{prop:restriction}
Under the hypotheses on $T$ above, let $f\circ T=f$, $\deg f\le3$, and suppose $f$ is bent. If $\dim U$ is even, then $f|_U$ is bent.
\end{proposition}
\begin{proof}
For each $a\in U\setminus\{0\}$, $D_af$ is quadratic, balanced, and $T$-invariant: $Ta=a$ implies $D_af(Tx)=D_af(x)$. Lemma~\ref{lem:reduction} shows that $(D_af)|_U=D_a(f|_U)$ is balanced. Apply the derivative characterization of bentness.
\end{proof}

\section{The restricted orbit space}
Write $n=tm$, where $t=2^s$, $s\ge1$, and $m$ is odd. The map $T=\sigma^t$ has order $m$, and its fixed subspace consists of repetitions of a block $y\in\F^t$. Under this identification $x_i=y_{i\bmod t}$.

\begin{lemma}\label{lem:orbit}
The restriction of any homogeneous cubic rotation-symmetric function belongs to the space $\mathcal F_t$ generated by all full-length cubic monomial orbit sums, the quadratic sums
\[Q_d(y)=\sum_{i=0}^{t-1}y_i y_{i+d},\quad 1\le d<t/2,\]
and $L(y)=\sum_{i=0}^{t-1}y_i$. Indices in these sums are modulo $t$.
\end{lemma}
\begin{proof}
Consider a cubic monomial with a three-element index set. Its stabilizer $H\le C_n$ acts freely on that set, so $h=|H|$ divides three and $h\in\{1,3\}$. Since $h\mid n=tm$ and $\gcd(h,t)=1$, $h\mid m$. Its orbit length is $n/h=t(m/h)$, with $m/h$ odd. Upon restriction, its orbit sum is $(m/h)$ times the sum of the $t$ cyclic translates of its reduced monomial. Since $m/h$ is odd, this is exactly that cyclic sum over $\F$. Contributions from distinct original cubic orbits then add in $\F$.

The reduced monomial has degree one, two or three, and no constant term. If its orbit length is $l\mid t$, its distinct translates occur $t/l$ times. When $l<t$, this multiplicity is even and the sum vanishes. All cubic monomial orbits in $t$ variables have length $t$, because their stabilizer divides both three and $t$. Quadratic orbits have length $t$ except for the antipodal orbit, of length $t/2$, which cancels. Equivalently, in the full-length case every antipodal monomial has total preimage multiplicity $2m$, hence even coefficient. A degree-one reduction gives $L$. This proves containment in $\mathcal F_t$; surjectivity of the restriction map is not required.
\end{proof}

\paragraph{Scope of the quadratic generators.}
There are $t/2$ distinct quadratic monomial orbits in the full rotation-symmetric space. Only $t/2-1$ have full length $t$. The remaining orbit is
\[
 P_t(y)=\sum_{i=0}^{t/2-1}y_i y_{i+t/2}.
\]
It is absent from $\mathcal F_t$: its reduced translates occur twice and cancel. Thus the divisibility assertion below is not an assertion about all rotation-symmetric functions of degree at most three. For example, $P_{16}$ is bent, with weight $32640\equiv128\pmod{256}$. Lower-degree terms arising from the restriction must be retained, but arbitrary lower-degree terms cannot be added to the stated space.

\begin{lemma}\label{lem:vanishing}
Every $g\in\mathcal F_t$ vanishes on $A_t=\{y:y_{i+t/2}=y_i\}$.
\end{lemma}
\begin{proof}
In a full orbit sum, pair the translate indexed by $i$ with that indexed by $i+t/2$. They evaluate equally on $A_t$ and cancel in characteristic two. The same pairing applies to the summands of $L$.
\end{proof}


\begin{lemma}[Removing the linear component]\label{lem:linear}
For a Boolean function $h$ and an affine form $\ell(x)=a\cdot x+c$,
\[W_{h+\ell}(b)=(-1)^c W_h(b+a).\]
Consequently, $h+\ell$ is bent if and only if $h$ is bent.
\end{lemma}
\begin{proof}
Substitute $h(x)+a\cdot x+c$ into the definition of the Walsh transform.
The map $b\mapsto b+a$ permutes its frequencies and the constant factor
does not change their absolute values.
\end{proof}
Define $\mathcal H_t$ to be the span of the full cubic and full quadratic
orbit sums, without $L$. Since every $g\in\mathcal F_t$ is $h+\epsilon L$
for $h\in\mathcal H_t$, excluding bent functions in $\mathcal H_t$ suffices.
This observation is an application of standard affine invariance; no
novelty is claimed for the identity itself.

\section{Global cubic nonexistence theorem}

We now state the structural result that supersedes the earlier dimension-by-dimension obstruction.

\begin{theorem}[Antipodal Rule]\label{thm:antipodal}
Let \(N=2^k\) with \(k\ge2\). If \(g:\F^N\to\F\) is rotation-symmetric, bent, and \(\deg g\le3\), then its ANF contains the antipodal quadratic orbit
\[
P_N(x)=\sum_{i=0}^{N/2-1}x_i x_{i+N/2}.
\]
\end{theorem}

\begin{proof}
Write \(N=2h\), \(d(u)=(u,u)\), and \(e(z)=(0,z)\). If the antipodal coefficient is zero, half-rotation pairing shows that, after adding a constant, \(g(d(u))=0\) for all \(u\). Put
\[
F_z(u)=g(u,u+z),\qquad A_z=\sum_u(-1)^{F_z(u)}.
\]
The invertible change \((u,z)\mapsto(u,u+z)\) preserves bentness. The Walsh transform of \(z\mapsto A_z\) is a Walsh slice of the transformed bent function; Parseval therefore gives
\[
\sum_z A_z^2=2^N.
\]
Since \(A_0=2^h\), every \(A_z\) with \(z\ne0\) vanishes. Thus \(F_{\mathbf1}\) is balanced.

Let \(J=\mathbf1^N\). Then \(D_Jg\) is balanced, rotation-symmetric and quadratic. A balanced rotation-symmetric quadratic in \(2^k\) variables has zero polar: if its polar \(B\) were nonzero, its radical \(R\) would be a proper rotation-invariant submodule of
\[
\F[X]/((X+1)^N),
\]
hence \(R\subseteq\operatorname{im}(S+I)\). For \(r=Sy+y\in R\), rotation invariance and alternation give \(Q(r)=B(r+y,y)=0\). The quadratic character identity would then give a nonzero zero-frequency Walsh sum, contradicting balance. Hence \(B_{D_Jg}=0\).

For \(\deg g\le3\), the third difference
\[
\Delta_3(p,q,r)=D_pD_qD_rg(0)
\]
is independent of the base point and is a symmetric trilinear form. Since \(g\circ d=0\),
\[
B_{F_{\mathbf1}}(u,v)
=\Delta_3(d(u),d(v),e(\mathbf1)).
\]
Moreover,
\[
B_{D_Jg}(d(u),e(v))=\Delta_3(d(u),e(v),J).
\]
Write \(e'(z)=(z,0)\), so \(J=e'(\mathbf1)+e(\mathbf1)\). The half-rotation fixes \(d(u)\) and exchanges \(e(v)\leftrightarrow e'(v)\) and \(e(\mathbf1)\leftrightarrow e'(\mathbf1)\). By invariance and trilinearity of \(\Delta_3\),
\[
\begin{aligned}
\Delta_3(d(u),e(v),J)
&=\Delta_3(d(u),e(v),e'(\mathbf1))+\Delta_3(d(u),e(v),e(\mathbf1))\\
&=\Delta_3(d(u),e'(v),e(\mathbf1))+\Delta_3(d(u),e(v),e(\mathbf1))\\
&=\Delta_3(d(u),d(v),e(\mathbf1)).
\end{aligned}
\]
Thus
\[
B_{F_{\mathbf1}}(u,v)=B_{D_Jg}(d(u),e(v))=0.
\]
A Boolean function with identically zero polar has no ANF terms of degree at least two, so \(F_{\mathbf1}\) is affine.

With the convention \((\sigma x)_i=x_{i-1\pmod N}\), write \(u=(u_0,\ldots,u_{h-1})\). A direct coordinate calculation gives
\[
\sigma(u,u+\mathbf1)
=(Ru+e_0,\;Ru+e_0+\mathbf1),
\]
where \(R\) is cyclic rotation on \(\F^h\) and \(e_0=(1,0,\ldots,0)\) (with the opposite rotation convention the unit vector occurs at the opposite endpoint). Hence \(F_{\mathbf1}(Ru+e_0)=F_{\mathbf1}(u)\). If
\(F_{\mathbf1}(u)=L(u)+c\), then \(L\circ R=L\), so
\(L(u)=\gamma\sum_i u_i\); the translation by \(e_0\) forces
\(\gamma=0\). Therefore \(F_{\mathbf1}\) is constant, contradicting its balance. The antipodal coefficient must be one.
\end{proof}

\begin{theorem}[Main theorem]\label{thm:global}
For every positive even integer \(n\), there is no homogeneous rotation-symmetric bent Boolean function of degree three in \(n\) variables.
\end{theorem}

\begin{proof}
Write \(n=2^s m\), where \(m\) is odd, and put \(t=2^s\). Let \(T=\sigma^t\). By Proposition~\ref{prop:restriction}, a hypothetical cubic homogeneous rotation-symmetric bent function \(f\) restricts to a bent rotation-symmetric function \(g=f|_{\Fix(T)}\) of degree at most three in \(t\) variables.

Lemma~\ref{lem:orbit} shows more specifically that the antipodal quadratic orbit is absent from \(g\): it has orbit length \(t/2\), so its \(t\) reduced translates occur twice and cancel. Equivalently, in a full cubic orbit each antipodal monomial receives \(2m\) preimages. Short cubic orbits have stabilizer three; since \(3\mid m\), their three support indices coincide modulo \(t\) and they reduce only to the total linear form.

If \(s\ge2\), Theorem~\ref{thm:antipodal} requires the antipodal orbit to be present in \(g\), a contradiction. If \(s=1\), then \(t=2\); folding leaves only degree at most one because the sole quadratic orbit is antipodal and cancels. Such a function cannot be bent in two variables. This proves the theorem.
\end{proof}

\paragraph{Relation to Sun--Shi--Liu--Fu (2026).}
Sun, Shi, Liu and Fu~\cite{sun} obtain additional conditional nonexistence results, including cubic criteria based on balanced derivatives, SANF/orbit structure and arithmetic hypotheses on \(n\). The proof above shares the derivative viewpoint but removes the odd part of \(n\) through the fixed-space reduction and then uses antipodal rigidity in the remaining power-of-two dimension. The comparison describes mechanism and scope; it is not an exhaustive priority claim.

\paragraph{Computational status.}
The earlier exact \(t=16\) divisible-code proof and the independent \(n=32\) certificates remain useful validation and historical motivation, but they are no longer logical premises of Theorem~\ref{thm:global}.

\section{Computational validation and supplementary material}
The proof of Theorem~\ref{thm:global} is entirely algebraic. Exact computations in dimensions 16 and 32, exhaustive low-dimensional checks, folding tests for odd multipliers, and nonhomogeneous cubic rotation-symmetric bent positive controls are retained in the repository as independent validation and historical documentation. None is a premise of the main theorem.

\section{Discussion}
The proof separates the cubic conjecture into two mechanisms. The odd-order fixed-space reduction removes the odd part of the dimension while preserving bentness because derivatives of cubic functions are quadratic. Cubic orbit folding then forbids the antipodal quadratic orbit in the restricted homogeneous space, whereas the Antipodal Rule forces precisely that orbit in every low-degree rotation-symmetric bent function in power-of-two dimension.

Homogeneity is essential to the folding obstruction, not to the Antipodal Rule. This is consistent with known nonhomogeneous cubic rotation-symmetric bent functions. For degree at least four, derivatives need not be quadratic, so the present fixed-space argument does not extend directly.

\begin{thebibliography}{10}
\bibitem{stanica} P.~St\u{a}nic\u{a} and S.~Maitra, Rotation symmetric Boolean functions---count and cryptographic properties, \emph{Discrete Applied Mathematics} 156(10) (2008), 1567--1580. \href{https://doi.org/10.1016/j.dam.2007.04.029}{doi:10.1016/j.dam.2007.04.029}.
\bibitem{meng} Q.~Meng, L.~Chen and F.-W.~Fu, On homogeneous rotation symmetric bent functions, \emph{Discrete Applied Mathematics} 158(10) (2010), 1111--1117. \href{https://doi.org/10.1016/j.dam.2010.02.009}{doi:10.1016/j.dam.2010.02.009}.
\bibitem{zhang} X.~Zhang and G.~Gao, On the conjecture about the nonexistence of rotation symmetric bent functions (2013). \href{https://arxiv.org/abs/1303.2282}{arXiv:1303.2282}.
\bibitem{cusick} T.~W.~Cusick and E.~M.~Sanger, Rotation symmetric bent Boolean functions for $n=2p$ (2017). \href{https://arxiv.org/abs/1708.09313}{arXiv:1708.09313}.
\bibitem{sun} L.~Sun, Z.~Shi, J.~Liu and F.-W.~Fu, On the conjecture about the nonexistence of homogeneous rotation symmetric bent functions, \emph{Designs, Codes and Cryptography} 94, article 93 (2026). \href{https://doi.org/10.1007/s10623-026-01848-4}{doi:10.1007/s10623-026-01848-4}.
\bibitem{rothaus} O.~S.~Rothaus, On ``bent'' functions, \emph{Journal of Combinatorial Theory, Series A} 20(3) (1976), 300--305. \href{https://doi.org/10.1016/0097-3165(76)90024-8}{doi:10.1016/0097-3165(76)90024-8}.
\bibitem{ward} H.~N.~Ward, Weight polarization and divisibility, \emph{Discrete Mathematics} 83(2--3) (1990), 315--326. \href{https://doi.org/10.1016/0012-365X(90)90015-A}{doi:10.1016/0012-365X(90)90015-A}.
\bibitem{sutang} S.~Su and X.~Tang, On the Systematic Constructions of Rotation Symmetric Bent Functions with Any Possible Algebraic Degrees, IACR ePrint 2015/451. \href{https://eprint.iacr.org/2015/451.pdf}{Preprint consulted}.
\end{thebibliography}
\end{document}

